\section{Conclusion}\label{sec:conclusion}

This study transforms a GA-based conference precursor into a controlled benchmark for neural intraday trend classification. It preserves the predecessor's practical insight---a compact, Information-Gain-selected technical representation can support metaheuristic hyperparameter search---but replaces randomized validation and a single-optimizer design with two primary chronological holdout protocols, a separate exploratory temporal cross-validation analysis, five budget-matched optimizers, and predictive and economic assessment. The five-minute feature vector at bar $t$ is paired with the regime label at $t+1$, derived from an hourly peak-and-trough sequence detected on resampled five-minute price data.

The current evidence indicates configuration-dependent changes between the two primary fitness protocols, which differ in both metric and train/validation composition. MLP leads the predictive aggregates under weighted-accuracy fitness, while the aligned next-bar replay gives 1D-CNN the largest mean net profit under both protocols; RF has the least negative mean drawdown under MCC/$F_1$ fitness. These results reinforce the central methodological lesson: selection objective, reported metric, and economic utility must be treated as related but distinct quantities.

The manuscript deliberately avoids a confirmatory claim about a universal best fitness function because only three paired seeds are available across every cell. The hourly pivot labels are retrospective. One training target crosses into validation, two validation targets cross into the nominal test interval, four per-ticker terminal rows lack a next-row target, and the temporal-CV folds were not purged for the one-row shift. A purged rerun is needed before claiming full test isolation; multiple market periods and more seeds are also needed before stronger generalization claims. The present benchmark is exploratory evidence for this fixed chronological interval.
